\section*{الفصل \ref{ch:b2:fourier} --- متسلسلات فورييه}

\begin{solution}{exo:b2:fourier:1}
تقتل الفردية المعاملات $a_n$. ومن أجل $n \geq 1$:
\[
b_n = \frac{2}{\pi}\int_0^{\pi} \sin nt\,\dd t
= \frac{2}{\pi}\cdot\frac{1 - (-1)^n}{n}
= \begin{cases} \frac{4}{\pi n} & n \text{ فردي},\\ 0 & n
\text{ زوجي}. \end{cases}
\]
ومنه $S(f)(t) = \frac{4}{\pi}\sum_{k\geq0} \frac{\sin\bigl((2k+1)t
\bigr)}{2k+1}$. وديريكليه عند $t = \frac\pi2$ (وهي نقطة اتصال،
بالقيمة $1$): $\sin\bigl((2k+1)\frac\pi2\bigr) = (-1)^k$، فيعطي
\[
1 = \frac4\pi \sum_{k\geq0}\frac{(-1)^k}{2k+1}
\quad\Longrightarrow\quad
\sum_{k\geq0}\frac{(-1)^k}{2k+1} = \frac{\pi}{4}
\quad\text{(لايبنتز)}.
\]
وعند القفزة $t = 0$: يساوي مجموع المتسلسلة $0 = \frac{f(0^+) +
f(0^-)}{2}$، كما يقتضي ديريكليه (فكل حدّ ينعدم:
وهو متوافق).
\end{solution}

\begin{solution}{exo:b2:fourier:2}
تقتل الزوجية المعاملات $b_n$؛ و$a_0 = \frac{1}{\pi}\int_{-\pi}^\pi\abs
t\,\dd t = \pi$، ومن أجل $n \geq 1$:
\[
a_n = \frac{2}{\pi}\int_0^\pi t\cos nt\,\dd t
= \frac{2}{\pi}\cdot\frac{(-1)^n - 1}{n^2}
= \begin{cases} -\frac{4}{\pi n^2} & n \text{ فردي},\\ 0 & n
\text{ زوجي}, \end{cases}
\]
(بمكاملة واحدة بالتجزئة). ومنه
\[
\abs t = \frac{\pi}{2} - \frac{4}{\pi}\sum_{k\geq0}
\frac{\cos\bigl((2k+1)t\bigr)}{(2k+1)^2} ,
\]
بتقارب \emph{ناظمي} (لأن $\sum (2k+1)^{-2} < \infty$) --- كما
يتنبأ البند (1) في \cref{thm:b2:fourier:parseval} من أجل هذه الدالة المتصلة
من الصنف $C^1$ على قطع. وعند $t = 0$:
\[
0 = \frac\pi2 - \frac4\pi\sum_{k\geq0}\frac{1}{(2k+1)^2}
\quad\Longrightarrow\quad
\sum_{k\geq0}\frac{1}{(2k+1)^2} = \frac{\pi^2}{8} .
\]
وبفصل $\sum \frac{1}{n^2}$ إلى جزأين فردي وزوجي: $S =
\frac{\pi^2}{8} + \frac S4$، ومنه $S = \frac{\pi^2}{6}$: أي بازل من جديد.
\end{solution}

\begin{solution}{exo:b2:fourier:3}
الزوجية: $b_n = 0$؛ و$a_0 = \frac{1}{\pi}\int_{-\pi}^{\pi} t^2 =
\frac{2\pi^2}{3}$؛ وتعطي مكاملتان بالتجزئة أن $a_n =
\frac{4(-1)^n}{n^2}$ (مع $n \geq 1$). ومنه
\[
t^2 = \frac{\pi^2}{3} + 4\sum_{n\geq1}
\frac{(-1)^n}{n^2}\cos nt
\qquad (\abs t \leq \pi),
\]
وهي متقاربة \omterm{def:b2:funcseq:series}{ناظميًا}. وعند $t = 0$: $0 = \frac{\pi^2}{3} +
4\sum\frac{(-1)^n}{n^2}$، أي $\sum_{n\geq1}
\frac{(-1)^{n+1}}{n^2} = \frac{\pi^2}{12}$. وبارسفال:
\[
\frac{1}{2\pi}\int_{-\pi}^{\pi} t^4\,\dd t = \frac{\pi^4}{5}
= \frac{a_0^2}{4} + \frac12\sum_{n\geq1} a_n^2
= \frac{\pi^4}{9} + 8\sum_{n\geq1}\frac{1}{n^4} ,
\]
ومنه $\sum \frac{1}{n^4} = \frac18\bigl(\frac{\pi^4}{5} -
\frac{\pi^4}{9}\bigr) = \frac{\pi^4}{90}$.
\end{solution}

\begin{solution}{exo:b2:fourier:4}
إذا كانت $f$ زيادةً على ذلك من الصنف $C^1$ على قطع: فيعطي بارسفال (المبرهَن)
أن $\norm f_2^2 = \sum\abs{c_n}^2 = 0$، وتفرض الإيجابية التامة لتكامل
الدالة المتصلة $\abs f^2$ أن $f = 0$.

ومن أجل $f$ المتصلة فحسب، نقبل بارسفال العام: فالبرهان نفسه في سطر
واحد. (والخطوط العريضة لطريق الكثافة: يبيّن التقريب المثلثي من
نمط فييير أو فايرشتراس أن كثيرات الحدود المثلثية
كثيفة \omterm{def:b2:nvs:norm}{بالمعيار} $\norm\cdot_2$ بين الدوال الدورية المتصلة؛ وبما أن $f
\perp$ كلها، فإن $\norm f_2^2 = \langle f, f - P\rangle \leq
\norm f_2\norm{f - P}_2$ من أجل المقرِّبات $P$، فيُفرض $\norm f_2 =
0$.)
\end{solution}

\begin{solution}{exo:b2:fourier:5}
الزوجية: $b_n = 0$؛
\[
a_n = \frac{2}{\pi}\int_0^\pi \cos(\alpha t)\cos(nt)\,\dd t
= \frac{2}{\pi}\cdot
\frac{(-1)^n\,\alpha\sin(\pi\alpha)}{\alpha^2 - n^2}
\]
(بتحويل الجداء إلى مجموع، ثم المكاملة؛ مع $a_0 =
\frac{2\sin(\pi\alpha)}{\pi\alpha}$). وديريكليه عند $t = \pi$ (وهي
نقطة اتصال للتمديد الدوري، وتتفق قيمتاها من جهة واحدة
بالزوجية):
\[
\cos(\pi\alpha)
= \frac{\sin(\pi\alpha)}{\pi\alpha}
+ \sum_{n\geq1} \frac{2\alpha\sin(\pi\alpha)}{\pi(\alpha^2 -
n^2)}\,(-1)^n\cos(n\pi)
= \frac{\sin(\pi\alpha)}{\pi}\Bigl(\frac{1}{\alpha} +
\sum_{n\geq1}\frac{2\alpha}{\alpha^2 - n^2}\Bigr),
\]
باستعمال $(-1)^n\cos n\pi = 1$. والقسمة على $\sin(\pi\alpha)/\pi$
تعطي نشر ظل التمام.
\end{solution}

\begin{solution}{exo:b2:fourier:6}
بتكرار $c_n(f') = \iu n\,c_n(f)$ ($k$ مرة، بالمكاملة بالتجزئة
عبر القطع من الصنف $C^{k}$ بقيم حدّية متوافقة):
$c_n(f^{(k)}) = (\iu n)^k c_n(f)$. ومعاملات الدالة
$f^{(k)}$ المتصلة على قطع محدودة (بل $\to 0$، ببيسل):
\[
\abs{c_n(f)} = \frac{\abs{c_n(f^{(k)})}}{\abs n^k}
= O\bigl(\abs n^{-k}\bigr) .
\]
\end{solution}

\begin{solution}{exo:b2:fourier:7}
بارسفال من أجل $f$ ومن أجل $f'$ (وكلاهما مشروع: فالدالة $f$ من الصنف $C^1$، و$f'$
متصلة على قطع --- بل متصلة):
\[
\frac{1}{2\pi}\int \abs f^2 = \sum_{n\neq0} \abs{c_n}^2
\quad (c_0 = 0 \text{ بفرضية انعدام المتوسط}),
\qquad
\frac{1}{2\pi}\int \abs{f'}^2 = \sum_{n\neq0} n^2\abs{c_n}^2 .
\]
وحدًّا حدًّا، $n^2\abs{c_n}^2 \geq \abs{c_n}^2$ من أجل $\abs n \geq 1$:
فتنتج المتراجحة. وتفرض المساواة أن $(n^2 - 1)\abs{c_n}^2 = 0$
من أجل كل $n$، أي $c_n = 0$ من أجل $\abs n \geq 2$: أي $f(t) =
c_1\eu^{\iu t} + c_{-1}\eu^{-\iu t} = a\cos t + b\sin t$ (في الصورة
الحقيقية)؛ وعكسيًا تعطي هذه الدوال $f$ المساواة.
\end{solution}

\begin{solution}{exo:b2:fourier:8}
من \cref{exo:b2:fourier:1}، $S_{2N-1}(x) =
\frac4\pi\sum_{k=0}^{N-1}\frac{\sin((2k+1)x)}{2k+1}$. وعند $x_N =
\frac{\pi}{2N}$، مع $u_k = (2k+1)x_N = \frac{(2k+1)\pi}{2N}$:
\[
S_{2N-1}(x_N) = \frac{4}{\pi}\sum_{k=0}^{N-1}\frac{\sin u_k}{2k+1}
= \frac{4}{\pi}\sum_{k=0}^{N-1}\frac{\sin u_k}{u_k}\cdot
\frac{u_k}{2k+1}
= \frac{2}{\pi}\sum_{k=0}^{N-1}\frac{\sin u_k}{u_k}\cdot
\frac{\pi}{N},
\]
لأن $\frac{u_k}{2k+1} = \frac{\pi}{2N}$. والنقاط $u_k$ هي
النقاط الوسطى للفترات الجزئية $N$ من $\intcc{0}{\pi}$ ذوات الطول
$\frac{\pi}{N}$: ومن ثم فالمجموع مجموع ريمان بالنقاط الوسطى للدالة
المتصلة $u \mapsto \frac{\sin u}{u}$، ومن ثم يتقارب إلى
\[
\frac{2}{\pi}\int_0^\pi \frac{\sin u}{u}\,\dd u
\approx \frac{2}{\pi}\times 1.8519 \approx 1.179 .
\]
فتتجاوز المجاميع الجزئية قرب القفزة القيمةَ $1$ بنسبة $\approx
18\%$ من نصف القفزة إلى الأبد: أي ظاهرة غيبس، مصوغةً كمّيًا.
\end{solution}

\begin{solution}{exo:b2:fourier:9}
من $\cos 3t = 4\cos^3t - 3\cos t$:
\[
\cos^3 t = \frac{3\cos t + \cos 3t}{4},
\qquad
\sin^2t\,\cos t = \cos t - \cos^3 t
= \frac{\cos t - \cos 3t}{4} .
\]
وكلٌّ منهما كثير حدود مثلثي، ومن ثم فهو متسلسلة فورييه
الخاصة به (بوحدانية المعاملات: إذ إن نشرين
سيختلفان بكثير حدود مثلثي كل
معاملاته معدومة). ومن أجل $\cos^3t$: $a_1 = \frac34$، $a_3 =
\frac14$، وكل المعاملات $a_n$ الأخرى وكل $b_n$ معدومة؛ و$c_{\pm1} =
\frac38$، $c_{\pm3} = \frac18$. ومن أجل $\sin^2t\cos t$: $a_1 =
\frac14$، $a_3 = -\frac14$؛ و$c_{\pm1} = \frac18$، $c_{\pm3} =
-\frac18$.
\end{solution}

\begin{solution}{exo:b2:fourier:10}
بالحساب المباشر:
\[
c_n = \frac{1}{2\pi}\int_{-\pi}^{\pi}\eu^{(a - \iu n)t}\dd t
= \frac{\eu^{(a-\iu n)\pi} - \eu^{-(a - \iu n)\pi}}
{2\pi(a - \iu n)}
= \frac{(-1)^n\sinh(a\pi)}{\pi(a - \iu n)} ,
\]
باستعمال $\eu^{\pm\iu n\pi} = (-1)^n$. وعند $t = \pi$ يقفز التمديد
الدوري من $\eu^{a\pi}$ إلى $\eu^{-a\pi}$؛ ويعطي ديريكليه
(بالمجاميع الجزئية المتناظرة)
\[
\cosh(a\pi) = \sum_{n\in\Z}(-1)^n c_n\,
= \frac{\sinh(a\pi)}{\pi}\Bigl(\frac1a +
\sum_{n\geq1}\Bigl(\frac{1}{a - \iu n} +
\frac{1}{a + \iu n}\Bigr)\Bigr)
= \frac{\sinh(a\pi)}{\pi}\Bigl(\frac1a +
\sum_{n\geq1}\frac{2a}{a^2 + n^2}\Bigr) ,
\]
مع تلاشي الأجزاء التخيلية للحدود المتزاوجة. ثم نقسم على
$\sinh(a\pi)$:
\[
\coth(\pi a) = \frac{1}{\pi a} +
\sum_{n\geq1}\frac{2a}{\pi(a^2 + n^2)} ,
\]
وهو توأم \cref{exo:b2:fourier:5} الزائدي.
\end{solution}

\begin{solution}{exo:b2:fourier:11}
التبادلية: نعوّض $s = x - t$ ونستعمل دورية
الدالة المكامَلة. ومن أجل $c_n(f * g)$، تكون الدالة المكامَلة $(x, t) \mapsto
f(x-t)g(t)\eu^{-\iu nx}$ متصلة؛ ويتفق التكاملان
المتكرران (فكلاهما، بوصفه دالة في الحدّ الأعلى للمتغير
الخارجي، ينعدم عند الطرف الأيسر وله المشتق نفسه ---
فالاتصال مع
\cref{thm:b2:integration:continuity} يبرّران اشتقاق
التكامل المتكرر). ومنه
\[
c_n(f*g) = \frac{1}{2\pi}\int_{-\pi}^{\pi} g(t)\,\eu^{-\iu nt}
\Bigl(\frac{1}{2\pi}\int_{-\pi}^{\pi}
f(x-t)\,\eu^{-\iu n(x-t)}\dd x\Bigr)\dd t
= c_n(f)\,c_n(g),
\]
لأن التكامل الداخلي يساوي $c_n(f)$ من أجل كل $t$ (بالتعويض
والدورية). وأخيرًا تقول \cref{lem:b2:fourier:kernel} إن
$S_N(f)(x) = \frac{1}{2\pi}\int f(x+u)D_N(u)\dd u$؛ ويحوّل
التعويض $u \mapsto -t$ وزوجيةُ $D_N$ هذا
إلى $(f * D_N)(x)$.
\end{solution}

\begin{solution}{exo:b2:fourier:12}
\begin{enumerate}
  \item مع $w = r\eu^{\iu t}$:
        \[
        P_r(t) = 1 + 2\,\Re\frac{w}{1 - w}
        = \Re\frac{1 + w}{1 - w}
        = \frac{1 - \abs w^2}{\abs{1 - w}^2}
        = \frac{1 - r^2}{1 - 2r\cos t + r^2} > 0 .
        \]
        والمتوسط $1$: فالمكاملة حدًّا حدًّا للمتسلسلة المتقاربة
        \omterm{def:b2:funcseq:series}{ناظميًا} لا تُبقي إلا $n = 0$.
  \item من أجل $\delta \leq \abs t \leq \pi$: $\cos t \leq
        \cos\delta$، ومنه $P_r(t) \leq \frac{1 - r^2}{1 -
        2r\cos\delta + r^2}$، ويؤول مقامه إلى $2 -
        2\cos\delta > 0$ في حين أن بسطه يؤول إلى $0$:
        أي \omterm{def:b2:funcseq:def}{التقارب المنتظم} إلى $0$ خارج أي جوار للنقطة
        $0$.
  \item والمكاملة حدًّا حدًّا (\omterm{def:b2:funcseq:series}{بالتقارب الناظمي} في $t$)
        تعطي $(f * P_r)(x) = \sum_n r^{\abs
        n}c_n(f)\eu^{\iu nx}$. وحجة المتطابقة التقريبية:
        فبمتوسط $1$ وبالإيجابية،
        \[
        \abs{(f*P_r)(x) - f(x)}
        \leq \frac{1}{2\pi}\int_{-\pi}^{\pi}
        \abs{f(x-t) - f(x)}\,P_r(t)\,\dd t ,
        \]
        نفصل عند $\abs t = \delta$: فلا يتجاوز $\varepsilon$
        (بهاينه) مضافًا إليه $2\norm f_\infty\sup_{\delta\leq\abs
        t\leq\pi}P_r \to \varepsilon$: أي \omterm{def:b2:funcseq:def}{التقارب المنتظم} حين
        $r \to 1^-$. وهذا هو \omterm{thm:b2:series:abel}{جمع آبل} لمتسلسلة
        فورييه --- أي توأم نظرية الحدّ في فصل متسلسلات
        القوى.
\end{enumerate}
\end{solution}

\begin{solution}{pb:b2:fourier:1}
\textbf{1.} أخذ متوسط \cref{lem:b2:fourier:kernel} على $n = 0,
\dots, N-1$ (بخطية التكامل) يعطي $\sigma_N(f)(x) =
\frac{1}{2\pi}\int f(x+u)F_N(u)\dd u$؛ ومتوسط كل $D_n$ هو
$1$، ومن ثم فمتوسط $F_N$ هو $1$.

\textbf{2.} مع $\eu^{\iu u} - 1 = 2\iu\,\eu^{\iu
u/2}\sin\frac u2$:
\[
\sum_{n=0}^{N-1}\sin\Bigl(\Bigl(n + \frac12\Bigr)u\Bigr)
= \Im\Bigl[\eu^{\iu u/2}\,\frac{\eu^{\iu Nu} - 1}{\eu^{\iu u}
- 1}\Bigr]
= \Re\,\frac{1 - \eu^{\iu Nu}}{2\sin\frac u2}
= \frac{1 - \cos Nu}{2\sin\frac u2}
= \frac{\sin^2\frac{Nu}2}{\sin\frac u2} .
\]
وبالقسمة على $N\sin\frac u2$:
\[
F_N(u) = \frac1N\sum_{n=0}^{N-1}
\frac{\sin\bigl((n+\frac12)u\bigr)}{\sin\frac u2}
= \frac{1}{N}\,
\frac{\sin^2\frac{Nu}{2}}{\sin^2\frac u2} \geq 0 .
\]

\textbf{3.} على $\delta \leq \abs u \leq \pi$: $\sin^2\frac u2
\geq \sin^2\frac\delta2$ و$\sin^2\frac{Nu}2 \leq 1$: ومنه $F_N
\leq \frac{1}{N\sin^2(\delta/2)} \to 0$، بانتظام هناك.

\textbf{4.} بكون المتوسط واحدًا، $\sigma_Nf(x) - f(x) =
\frac{1}{2\pi}\int\bigl(f(x+u) - f(x)\bigr)F_N(u)\dd u$. ومن أجل
$\varepsilon$ معطى، يعطي هاينه $\delta$ بحيث $\abs{f(x+u) - f(x)}
\leq \varepsilon$ من أجل $\abs u \leq \delta$، بانتظام في $x$.
وعندئذٍ، باستعمال $F_N \geq 0$ وكون متوسطها واحدًا،
\[
\abs{\sigma_Nf(x) - f(x)}
\leq \varepsilon +
2\norm f_\infty\cdot\frac{1}{2\pi}
\int_{\delta\leq\abs u\leq\pi}F_N
\leq \varepsilon + \frac{2\norm
f_\infty}{N\sin^2\frac\delta2}
\leq 2\varepsilon
\]
من أجل $N$ الكبيرة، وبانتظام في $x$: وهي مبرهنة فييير.

\textbf{5.} $F_N$ زوجية ومتوسطها $1$: فيحمل كل نصف
$\intcc{0}{\pi}$، $\intcc{-\pi}{0}$ المتوسطَ $\frac12$.
وعندئذٍ
\[
\sigma_Nf(x) - \frac{f(x^+)+f(x^-)}{2}
= \frac{1}{2\pi}\int_0^\pi\bigl(f(x+u) -
f(x^+)\bigr)F_N\,\dd u
+ \frac{1}{2\pi}\int_{-\pi}^0\bigl(f(x+u) -
f(x^-)\bigr)F_N\,\dd u ;
\]
على كل نصف، نفصل عند $\abs u = \delta$ حيث تكون النهاية من جهة
واحدة قريبة بمقدار $\varepsilon$، ويقتل السؤال 3 الجزء
البعيد: فيؤول التكاملان إلى $0$. والحدّ: يعطي $F_N \geq 0$
أن $\abs{\sigma_Nf(x)} \leq
\frac{1}{2\pi}\int\abs{f(x+u)}F_N \leq \norm f_\infty$.

\textbf{6.} كل $\sigma_N(f)$ كثير حدود مثلثي
(فهو متوسط للمقادير $S_n(f)$، مع $n < N$)، و$\sigma_N(f) \to f$
بانتظام: وهي مبرهنة فايرشتراس المثلثية.

\textbf{7.} يجعل $c_n(f) = 0$ من أجل كل $n$ كلَّ $S_n(f) = 0$،
ومن ثم كل $\sigma_N(f) = 0$؛ وبفييير، $f = \lim\sigma_Nf =
0$. وبتطبيق هذا على فرق: تتحدد الدوال الدورية
المتصلة \omterm{def:b2:fourier:coefficients}{بمعاملات فورييه} الخاصة بها.

\textbf{8.} $S_Nf$ هو الإسقاط المتعامد للدالة $f$ على
$\mathcal T_N$ (\cref{prop:b2:fourier:bessel})، ومن ثم فهو يصغّر
$\norm{f - P}_2$ على $P \in \mathcal T_N$؛ وبما أن $\sigma_Nf
\in \mathcal T_{N-1} \subseteq \mathcal T_N$:
\[
\norm{f - S_Nf}_2 \leq \norm{f - \sigma_Nf}_2
\leq \norm{f - \sigma_Nf}_\infty \to 0
\]
(والمتراجحة الوسطى لأن متوسط $\abs\cdot^2$ لا يتجاوز
مربّع النهاية العليا). ثم تمرّ مبرهنة فيثاغورس $\norm f_2^2 = \sum_{\abs
n\leq N}\abs{c_n}^2 + \norm{f - S_Nf}_2^2$ إلى
النهاية: أي بارسفال من أجل كل دالة $f$ دورية متصلة.

\textbf{9.} لتكن $t_1, \dots, t_p$ قفزات $f$ في
دور واحد، مع $M = \norm f_\infty$. ومن أجل $\eta$ الصغيرة، نعرّف $g = f$
خارج الفترات $\intoo{t_j - \eta}{t_j + \eta}$ وبالوتر
الأفيني عبر كل فترة كهذه: فتكون $g$ متصلة و
دورية، و$\norm g_\infty \leq M$، و
\[
\norm{f - g}_2^2 \leq \frac{1}{2\pi}\,p\cdot(2M)^2\cdot2\eta
\leq \varepsilon^2
\]
من أجل $\eta$ الصغيرة. ويجعل بيسل $S_N$ تقلّصًا من أجل
$\norm\cdot_2$، ومن ثم
\[
\norm{f - S_Nf}_2
\leq \norm{f - g}_2 + \norm{g - S_Ng}_2 + \norm{S_N(g -
f)}_2
\leq 2\varepsilon + \norm{g - S_Ng}_2 ,
\]
ويعطي السؤال 8 أن $\limsup_N\norm{f - S_Nf}_2 \leq
2\varepsilon$ من أجل كل $\varepsilon$: أي $\norm{f - S_Nf}_2 \to
0$، وتعطي مبرهنة فيثاغورس بارسفال من أجل كل دالة $f$
متصلة على قطع: فصار ``المقبول'' في
\cref{thm:b2:fourier:parseval} مبرهنةً الآن.

\textbf{10.} للموجة المربّعة $\norm f_\infty = 1$، ومن ثم
يعطي السؤال 5 أن $\abs{\sigma_N(f)} \leq 1$ في كل مكان ومن أجل
كل $N$ --- فلا تجاوز أبدًا --- في حين أن
\cref{exo:b2:fourier:8} يبيّن أن $\sup_xS_{2N-1}(f)(x) \to
\approx 1.179$. والخاصية المسؤولة عن ذلك: $F_N \geq 0$، ومن ثم فإن
$\sigma_Nf(x)$ \emph{متوسط} موزون لقيم $f$
ولا يمكن أن يخرج من $\intcc{\min f}{\max f}$ أبدًا؛ أما $D_N$ فتأخذ
قيمًا سالبة، ومن ثم يستطيع $S_N$ ذلك.

\textbf{11.} $\abs{\sin N\theta} \leq N\abs{\sin\theta}$ بالتراجع
($\abs{\sin(N{+}1)\theta} \leq
\abs{\sin N\theta}\abs{\cos\theta} +
\abs{\cos N\theta}\abs{\sin\theta} \leq
(N+1)\abs{\sin\theta}$): فمع $\theta = \frac u2$،
\[
F_N(u) = \frac{\sin^2\frac{Nu}2}{N\sin^2\frac u2}
\leq \frac{N^2\sin^2\frac u2}{N\sin^2\frac u2} = N .
\]
ويعطي تقعّر $\sin$ على $\intcc{0}{\frac\pi2}$
أن $\sin\frac u2 \geq \frac{u}{\pi}$ من أجل $0 \leq u \leq \pi$، ومنه
$F_N(u) \leq \frac{1}{N(u/\pi)^2} = \frac{\pi^2}{Nu^2}$.

\textbf{12.} بالزوجية وبالفصل عند $\frac1N$:
\[
\frac{1}{2\pi}\int_{-\pi}^{\pi}\abs uF_N
= \frac1\pi\int_0^\pi uF_N
\leq \frac1\pi\Bigl(\int_0^{1/N}uN\,\dd u +
\int_{1/N}^{\pi}\frac{\pi^2}{Nu}\,\dd u\Bigr)
= \frac{1}{2\pi N} + \frac{\pi\ln(\pi N)}{N} .
\]
ومن أجل $N \geq 2$: $\ln(\pi N) \leq \bigl(1 +
\frac{\ln\pi}{\ln2}\bigr)\ln N \leq 2.66\ln N$ و
$\frac{1}{2\pi N} \leq \frac{\ln N}{N}$، ومن ثم فالعزم
$\leq \frac{9\ln N}{N}$.

\textbf{13.} من أجل $f$ الليبشيتزية بالثابت $L$:
\[
\abs{\sigma_Nf(x) - f(x)}
\leq \frac{1}{2\pi}\int\abs{f(x+u) - f(x)}F_N(u)\dd u
\leq L\cdot\frac{1}{2\pi}\int\abs uF_N
\leq \frac{9L\ln N}{N},
\]
بانتظام في $x$.

\textbf{14.} يعطي $c_0(e_1) = 0$ أن $S_0(e_1) = 0$، في حين أن
$S_n(e_1) = e_1$ من أجل $n \geq 1$: ومنه $\sigma_N(e_1) =
\frac{N-1}{N}e_1$، أي $\norm{\sigma_Ne_1 - e_1}_\infty =
\frac1N$. وحتى من أجل هذه الإشارة التامة المحدودة النطاق يكون المعدل
$\frac1N$: فيتشبّع فييير، تمامًا كما يتشبّع مؤثر برنشتاين
عند $\frac1n$ (فورونوفسكايا، في مسألة نهاية الأسبوع
في فصل متتاليات الدوال).

\textbf{15.} إذا انعدمت $f$ على $\intoo{x-\delta}{x+\delta}$،
فإن $\sigma_Nf(x) = \frac{1}{2\pi}\int_{\delta \leq \abs u
\leq \pi}f(x+u)F_N(u)\dd u$، وقيمته المطلقة لا تتجاوز
$\frac{\norm f_\infty}{N\sin^2(\delta/2)} \to 0$: فمتوسطات تشيزارو
عند $x$ لا ترى إلا $f$ قرب $x$.

\textbf{16.} من أجل $\intcc ab$ و$\varepsilon$ معطيين، نختار
دالتين متصلتين $\varphi^\pm$ دوريتين بالدور $1$ وأفينيتين على قطع تحققان
$\varphi^- \leq \mathbf 1_{\intcc ab} \leq \varphi^+$ و
$\int_0^1(\varphi^+ - \varphi^-) \leq \varepsilon$ (بشبه منحرفات
ذوات ميول على فترات طولها الكلي $\varepsilon$).
عندئذٍ
\[
\limsup_N\frac{\#\{n \leq N : x_n \in \intcc ab\}}{N}
\leq \lim_N\frac1N\sum_{n\leq N}\varphi^+(x_n)
= \int_0^1\varphi^+ \leq b - a + \varepsilon ,
\]
وبالتناظر $\liminf \geq b - a - \varepsilon$: ومن ثم تؤول
النسبة إلى $b - a$.

\textbf{17.} من أجل $f = \eu^{2\iu\pi k\cdot}$ مع $k \neq 0$ تعطي
الفرضية النهايةَ $0 = \int_0^1f$؛ ومن أجل $k = 0$ يساوي الطرفان
$1$؛ وتعالج الخطية كل كثير حدود
مثلثي. ومن أجل $f$ المتصلة الدورية بالدور $1$ ومن أجل $\varepsilon >
0$، يعطي السؤال 6 (منقولًا بالمقدار $t = 2\pi x$)
كثير حدود مثلثي $P$ يحقق $\norm{f - P}_\infty \leq
\varepsilon$:
\[
\Bigl|\frac1N\sum_{n\leq N}f(x_n) - \int_0^1f\Bigr|
\leq 2\varepsilon +
\Bigl|\frac1N\sum_{n\leq N}P(x_n) - \int_0^1P\Bigr|
\longrightarrow 2\varepsilon .
\]
ومع السؤال 16: تكون $(x_n)$ متوزعة بالتساوي.

\textbf{18.} من أجل $k \neq 0$ ومن أجل $\alpha$ الأصمّ، $w =
\eu^{2\iu\pi k\alpha} \neq 1$:
\[
\Bigl|\sum_{n=1}^{N}w^n\Bigr|
= \Bigl|\frac{w(w^N - 1)}{w - 1}\Bigr|
\leq \frac{2}{\abs{1 - w}} ,
\]
وهو حدّ لا يتعلق بالمقدار $N$؛ والقسمة على $N$ تعطي محك
فايل، ويختم السؤال 17: أي إن $(\{n\alpha\})$
متوزعة بالتساوي.

\textbf{19.} تتلقّى كل فترة جزئية نسبة مقاربة
تساوي طوله، وعلى الخصوص عددًا غير منته من النقاط:
فالمتتالية $(\{n\alpha\})$ كثيفة. والتوزع المتساوي أقوى: فقد تكون
متتالية كثيفة وهي تقضي جلّ وقتها في زاوية
واحدة (فالكثافة تقول \emph{أين} تذهب المتتالية،
والتوزع المتساوي يقول \emph{كم مرة}).

\textbf{20.} يكون الرقم الرائد للمقدار $2^n$ هو $1$ إذا وفقط إذا كان $10^m \leq 2^n <
2\cdot10^m$ من أجل $m$ ما، أي إذا وفقط إذا $\{n\log_{10}2\} \in
\intco{0}{\log_{10}2}$. والصمم: فالمساواة $\log_{10}2 =
\frac pq$ ستعطي $2^q = 10^p = 2^p5^p$، وهذا مستحيل من أجل $p
\geq 1$ بوحدانية التفكيك. وتعطي مبرهنة فايل (السؤال 18
مع $\alpha = \log_{10}2$؛ وتُحصر الفترة نصف \omterm{def:b2:metric:topology}{المفتوحة}
بين فترتين مغلقتين بطولين قريبين) النسبةَ
$\log_{10}2 \approx 0.301$: فأرقام $2^n$ الأولى تتبع
قانون بنفورد.

\textbf{21.} $z$ من الصنف $C^1$، ومن ثم تتقارب متسلسلة فورييه الخاصة بها
\omterm{def:b2:funcseq:series}{ناظميًا} بمجموع $z$ (البند (1) في \cref{thm:b2:fourier:parseval})،
و$c_n(z') = \iu n\,c_n$. وبما أن $\abs{z'} = \frac{L}{2\pi}$
ثابتة،
\[
\int_0^{2\pi}\abs{z'}^2\dd t
= 2\pi\Bigl(\frac{L}{2\pi}\Bigr)^{\!2}
= \frac{L^2}{2\pi} ,
\]
ويعطي بارسفال مطبَّقًا على الدالة المتصلة $z'$
أن $\frac{1}{2\pi}\int_0^{2\pi}\abs{z'}^2 =
\sum_n\abs{\iu nc_n}^2$، أي $\frac{L^2}{2\pi} = 2\pi\sum_n
n^2\abs{c_n}^2$.

\textbf{22.} تنتج صيغة بارسفال المستقطبة $\frac{1}{2\pi}\int\conj
fg = \sum\conj{c_n(f)}c_n(g)$ من بارسفال مطبَّقًا على
$f + g$ و$f + \iu g$ (\omterm{def:b2:quadratic:def}{بمتطابقة الاستقطاب})، وكلتاهما
متصلة. ومع $f = z$، $g = z'$:
\[
A = \frac12\,\Im\int_0^{2\pi}\conj z\,z'
= \pi\,\Im\sum_n\conj{c_n}(\iu n c_n)
= \pi\sum_n n\abs{c_n}^2 .
\]

\textbf{23.} بجمع السؤالين 21--22:
\[
L^2 - 4\pi A
= 4\pi^2\sum_n n^2\abs{c_n}^2 -
4\pi^2\sum_n n\abs{c_n}^2
= 4\pi^2\sum_{n\in\Z}(n^2 - n)\abs{c_n}^2 \geq 0 ,
\]
لأن $n^2 - n = n(n-1) \geq 0$ من أجل كل عدد صحيح. وتفرض المساواة
أن $c_n = 0$ من أجل كل $n \notin \{0, 1\}$: أي $z(t) = c_0 +
c_1\eu^{\iu t}$، وهي دائرة مركزها $c_0$ ونصف قطرها
$\abs{c_1} = \frac{L}{2\pi}$ (بسرعة ثابتة). وهذا هو برهان هورفيتس
لمتراجحة القياس المتساوي: $A \leq \frac{L^2}{4\pi}$،
والدائرة وحدها.

\textbf{24.} الدائرة ذات نصف القطر $R$: $L = 2\pi R$، $A = \pi
R^2$: أي $L^2 = 4\pi^2R^2 = 4\pi A$: فالمساواة. والمربّع ذو الضلع
$a$: $L^2 = 16a^2 > 4\pi a^2 = 4\pi A$ (لأن $16 > 4\pi \approx
12.57$). ومن أجل $n$ السالبة، يكون $n^2 - n = n(n - 1)$ جداء
عددين صحيحين سالبين: فهو موجب --- ومن ثم تكلّف الأنماط
ذات اللفّ العكسي المساحةَ مرتين. ودخلت السرعة الثابتة في السؤال 21،
فحوّلت $\int\abs{z'}^2$ إلى $\frac{L^2}{2\pi}$؛ ومن أجل
سرعة غير ثابتة، تعطي كوشي--شوارتز أن $\int\abs{z'}^2 \geq
\frac{(\int\abs{z'})^2}{2\pi} = \frac{L^2}{2\pi}$، ومن ثم
تصمد المتراجحة، والدائرة ما تزال حالة المساواة الوحيدة.

\textbf{25.} (1) كل شيء يتدفق من $F_N \geq 0$ (بمتوسط
واحدي وبالتركّز)؛ ولنواة $D_N$ متوسط واحدي وتركّز
للتذبذب لكن بلا إيجابية، وغيبس هو الثمن.
(2) وتستلزم القابلية للجمع بمعنى تشيزارو لمتسلسلة فورييه
قابليتَها للجمع بمعنى آبل بالمجموع نفسه (فروبينيوس، المبرهَن في
مسألة نهاية الأسبوع في فصل متسلسلات القوى) --- وطريقُ نواة
بواسون في \cref{exo:b2:fourier:12} هو بالضبط طريقة آبل.
(3) واحتاجت مبرهنة فايل التقريبَ المنتظم وحده
(السؤال 6)؛ واحتاجت متراجحة القياس المتساوي بارسفال
نفسه (الأسئلة 8 و21--22). (4) ويبرهن مجلد السنة الثالثة على
التمام: فالأسّيات تشكّل أساسًا هيلبرتيًا للفضاء $L^2$،
ويصير بارسفال تقايسًا بين فضاءي هيلبرت، وتصير مبرهنة
فييير القولَ بأن هذا التقايس قابل للحساب
بمتوسطات موجبة.
\end{solution}
